% !TeX root = ../../main.tex
\subsection{Intersection conventions and normalization}
\label{subsec:s02:01}

We work in the rational Chow ring
$A^*(\MM_{g,n})_{\mathbb Q}$.  Degree always means complex codimension:
for each integer $a\geq0$, a class in
$A^a(\MM_{g,n})_{\mathbb Q}$ has degree $a$.  Integration is the degree of a
zero-cycle on the stack, including its stabilizer weights.  For a class
with several homogeneous components, the integral selects the component
of degree $3g-3+n$.  We allow extension of the coefficient field when
scalar parameters occur.

We set $\lambda_0=1$ and $\lambda_j=0$ for every integer $j>g$.

If $\theta$ is a class with only positive-degree components, its
exponential means the formal sum
$e^\theta=\sum_{k\geq0}\theta^k/k!$, where $k$ ranges over the
nonnegative integers.  On $\MM_{g,n}$ this sum is truncated in degrees
above $3g-3+n$.  The following elementary identity will restore the
factorials when we return from exponential integrals to powers.

\begin{lmm}[Exponentials and powers]
\label{lmm:s02:exp-to-power}
Let $w$ be an integer with $0\leq w\leq3g-3+n$, let
$Q\in A^w(\MM_{g,n})_{\mathbb Q}$ be homogeneous, and let
$L\in A^1(\MM_{g,n})_{\mathbb Q}$ be a divisor class.  For every scalar
$c$,
\begin{equation}
\label{eq:s02:exp-to-power}
\int_{\MM_{g,n}}Qe^{cL}
=\frac{c^{3g-3+n-w}}{(3g-3+n-w)!}
  \int_{\MM_{g,n}}QL^{3g-3+n-w}.
\end{equation}
\end{lmm}

\begin{prf}
In the exponential expansion the term $Qc^kL^k/k!$ has degree $w+k$.
The integral therefore retains exactly the term with
$k=3g-3+n-w$.
\end{prf}

For example, if $g\geq2$ and $h$ is an integer with $0\leq h\leq g$,
\begin{equation}
\label{eq:s02:hodge-factorials}
A_g=(3g-3)!\int_{\MM_g}e^{\lambda_1},
\qquad
B_g(h)=(g_2+h)!\int_{\MM_g}\lambda_{g-h}e^{\lambda_1}.
\end{equation}
These are exact identities at finite genus.

The parameter conventions for subsequent asymptotic statements are as
follows.
\begin{center}
\begin{tabular}{ll}
Parameter & Convention as $g\to\infty$ \\
\hline
$j_1,\ldots,j_\ell$ and $\ell$ & fixed for low-index Hodge insertions, \\
$h$ & fixed, nonnegative, for $\lambda_{g-h}$, \\
expansion order & fixed for each stated expansion, \\
$n$ & fixed only when the statement so specifies.
\end{tabular}
\end{center}
Implicit constants in error terms may depend on these fixed data.
In graph sums, vertex valences and the number of nodes may grow with
$g$; summation of error terms will require estimates uniform over the
stated indexing set.  Formal truncation in a graded ring and such
uniform asymptotic estimates serve different purposes.
